\abstract % Структурный элемент: РЕФЕРАТ

\keywords{БЕСПИЛОТНЫЙ АВТОМОБИЛЬ, ЛИДАРНАЯ ОДОМЕТРИЯ, РЕГИСТРАЦИЯ ОБЛАКОВ ТОЧЕК, ICP, ПЕРЕКРЫТИЕ ПОЛЯ ЗРЕНИЯ, ФИЛЬТР КАЛМАНА, BOREAS-RT}

Объектом исследования в данной работе являются методы локализации и оценки движения беспилотного автомобиля на основе лидарных данных.

Цель работы --- разработка алгоритма лидарной одометрии беспилотного автомобиля, учитывающего перекрытие поля зрения и обеспечивающего устойчивую оценку движения на коротком участке пути.

Основное содержание работы состоит в рассмотрении особенностей технологии лидарной одометрии, анализе алгоритма ICP и ограничений регистрации облаков точек при динамическом перекрытии поля зрения, разработке устойчивого алгоритма лидарной одометрии с использованием ICP, оценки преобразования между парами точек с учётом весов, модели постоянной скорости и фильтра Калмана, а также в построении экспериментального конвейера для оценки качества предложенного подхода.

Основным результатом работы является устойчивый алгоритм лидарной одометрии, реализованный на языке Python, который включает предобработку облаков точек, регистрацию соседних лидарных кадров методом ICP, оценку преобразования между парами точек с учётом весов, дающих устойчивость к выбросам, и фильтр Калмана для сглаживания и предсказания скоростей.

Экспериментальное исследование проводится на данных набора Boreas-RT с использованием эталонных значений линейных и угловых скоростей. Качество работы алгоритма оценивается по метрикам среднего смещения, стандартного отклонения, MSE, MAE, RMSE и 95\%-квантиля абсолютной ошибки, а также по результатам сравнения различных конфигураций алгоритма.

Практическая значимость работы заключается в возможности использования предложенного подхода для повышения устойчивости лидарной одометрии беспилотного автомобиля в дорожных сценах, где крупные динамические объекты временно перекрывают значительную часть поля зрения лидара.